\section{Exact interlevel contact and the second-layer degree}\label{sec:second}

The first comparison supplies degree \(p^2\) uniformly, without
an Artin--Schreier coefficient calculation.

\begin{proposition}\label{prop:second}
For every \(e\ge2\), \(\alpha\in S_e\), and \(\beta\in S_1\),
\begin{equation}\label{eq:exact-contact}
 v(\alpha-\beta)=d_0:=\frac{2(p-1)^2}{p^2}.
\end{equation}
Moreover,
\begin{equation}\label{eq:pstep}
 [K(\alpha):K]\ge p^2,\qquad
 v(P^{\circ p}(\alpha)-\alpha)=2(p-1).
\end{equation}
In particular, \(M_2\) is irreducible of degree \(p^2\).
\end{proposition}

\begin{proof}
The exact first factorization is
\begin{equation}\label{eq:first-factor}
 P^{\circ p}(z)-z=z(z+s)M_1(z)V_1(z).
\end{equation}
Put \(\mu=(P^{\circ p})'(\beta)\).
Differentiating at \(\beta\) gives
\[
 \mu-1=\beta(\beta+s)M_1'(\beta)V_1(\beta).
\]
The first two factors have valuation \(r\), the last is a unit,
and each of the \(p-1\) root differences in \(M_1'(\beta)\)
has valuation \(2r\). Hence
\begin{equation}\label{eq:first-multiplier}
 v(\mu-1)=2r+(p-1)2r=2(p-1)<+\infty .
\end{equation}

For \(e\ge2\), both return polynomials defining \(Q_e\) vanish
at \(\beta\). Differentiate their polynomial quotient identity
before evaluation. By the chain rule and \eqref{eq:first-multiplier},
the derivative of the denominator is nonzero, and
\[
 Q_e(\beta)=
 \frac{\mu^{p^{e-1}}-1}{\mu^{p^{e-2}}-1}
 =(\mu-1)^{(p-1)p^{e-2}} .
\]
This is not substitution into an undefined \(0/0\) expression.
Since \(V_e(\beta)\) is a unit, it follows that
\begin{equation}\label{eq:first-average}
 \sum_{\alpha'\in S_e}v(\beta-\alpha')
 =2(p-1)^2p^{e-2}.
\end{equation}
All contacts are finite, because the ordinary periods differ.
Any two level-\(e\) points differ by a sum of one-step
displacements and thus have difference valuation at least \(2r\).
The average in \eqref{eq:first-average} is \(d_0<2r\).
Some contact is therefore below \(2r\); the unequal-valuation
triangle rule forces every contact to equal that one. Their
average is \(d_0\), proving \eqref{eq:exact-contact}.

Write \(L=K(\alpha)\), \(F=K(\beta)=L_1\), and \(B=LF\).
By Lemma~\ref{lem:rotation}, \([L:K]=p^h\), \(1\le h\le e\).
Suppose \(h=1\). Since \(p\) is odd, \(d_0\) has exact
denominator \(p^2\). The element \(\eta=\alpha-\beta\in B\)
therefore forces \([B:K]\ge p^2\); the opposite inequality
follows from the two degrees. Thus \(L\cap F=K\),
\(\Gal(B/K)\simeq C_p\times C_p\), and \(v_B=p^2v\).
Because \(h<e\), each nontrivial automorphism of \(L/K\)
has a rotation index divisible by \(p\), so moves \(\alpha\)
with valuation at least \(3r\). Every nontrivial automorphism
of \(F/K\) moves \(\beta\) with valuation \(2r\).
The restrictions can be chosen independently in the direct
product. Hence the minimum nonidentity displacement of
\(\eta\) has base valuation exactly \(2r\).

But \(v_B(\eta)=2(p-1)^2\) is positive and prime to \(p\).
Lemma~\ref{lem:prime-value} makes the first break of \(B/K\)
\[
 p^2(2r-d_0)=2(p-1).
\]
Its quotient \(F/K\) has break \(p-1\), contradicting upper
quotient compatibility. Consequently \(h\ge2\).
At \(e=2\), the orbit size also bounds the degree by \(p^2\),
so equality and irreducibility follow.
Finally evaluate \eqref{eq:first-factor} at \(\alpha\).
Both initial factors have valuation \(r\), the complementary
factor is a unit, and all \(p\) contacts equal \(d_0\). Thus
\[
 v(P^{\circ p}(\alpha)-\alpha)=2r+p d_0=2(p-1),
\]
which proves \eqref{eq:pstep}.
\end{proof}

The contact \eqref{eq:exact-contact} has denominator \(p^2\)
at every higher level. In particular, \(p^e d_0\) is divisible
by \(p\) when \(e\ge3\), so the prime-value argument cannot be
repeated with this contact to force degree \(p^e\).
The next section uses level-two clusters instead.
